% GENERATED FILE. Edit canonical lessons, then run npm run generate:books.
% canonical-source-hash: fnv1a64:5f1a9822ff8413de
% canonical-lessons: ES-C526-rescatar, ES-C526-rescate, ES-C526-salvar, ES-C526-sobrevivir, ES-C526-atrapar

\chapter{El rescate}
\label{ch:es-el-rescate}
\hlchaptermodality{526}

\begin{chapteropening}
\textbf{By the end of this chapter:} \emph{I can say that people were trapped, rescued and saved, and that they survived.}

Complete the last lesson of chapter 526: I can say that people were trapped, rescued and saved, and that they survived.
\end{chapteropening}

\section[rescatar]{rescatar — to rescue}

\label{lesson:ES-C526-rescatar}



\begin{quote}
Before the new one: say the Spanish for to help, then the Spanish for the injured person.
\end{quote}

\subsection*{You'll want to know: rescatar}
\textbf{rescatar} — \textquotedblleft{}to rescue\textquotedblright{}.

\begin{quote}
Rescataron a los niños del río.
\end{quote}

\begin{quote}
\emph{They rescued the children from the river.}
\end{quote}

\begin{cousinweb}
\emph{rescatar} comes from spoken Latin \emph{*reexcaptāre} 'to take back' (EtymDB), built on \emph{captāre} 'to try to catch', from \emph{capere} 'to take'.
\end{cousinweb}

\begin{grammarlens}[title={What you've built}]
One more word for telling, explaining and describing at B1.
\end{grammarlens}

\subsection*{Your turn}
\begin{itemize}
\raggedright
  \item \emph{Say it:} \emph{rescatar}
  \item \emph{Say it:} \emph{rescatar}, once more
  \item \emph{Say it:} a sentence of your own with \emph{rescatar}
  \item \emph{Recall:} say the Spanish for to help, then the Spanish for the injured person, then say \emph{rescatar} again
\end{itemize}

\begin{tcolorbox}[breakable,colback=teal!4,colframe=teal!35!black,title={Before you move on}]
What does rescatar mean? (\textquotedblleft{}To rescue\textquotedblright{}.) Say it once more.
\end{tcolorbox}
\section[el rescate]{el rescate — the rescue; the ransom}

\label{lesson:ES-C526-rescate}



\begin{quote}
Before the new one: say the Spanish for the injured person, then the Spanish for to rescue.
\end{quote}

\subsection*{You'll want to know: el rescate}
\textbf{el rescate} — \textquotedblleft{}the rescue; the ransom\textquotedblright{}.

\begin{quote}
El rescate duró toda la noche.
\end{quote}

\begin{quote}
\emph{The rescue lasted all night.}
\end{quote}

It is also the money asked for a hostage: \emph{pedir un rescate}.

\begin{cousinweb}
\emph{rescate} is built on \emph{rescatar} 'to rescue' (EtymDB), which goes back to spoken Latin \emph{*reexcaptāre}, from \emph{captāre} 'to try to catch'.
\end{cousinweb}

\begin{grammarlens}[title={What you've built}]
One more word for telling, explaining and describing at B1.
\end{grammarlens}

\subsection*{Your turn}
\begin{itemize}
\raggedright
  \item \emph{Say it:} \emph{el rescate}
  \item \emph{Say it:} \emph{el rescate}, once more
  \item \emph{Say it:} a sentence of your own with \emph{el rescate}
  \item \emph{Recall:} say the Spanish for the injured person, then the Spanish for to rescue, then say \emph{el rescate} again
\end{itemize}

\begin{tcolorbox}[breakable,colback=teal!4,colframe=teal!35!black,title={Before you move on}]
What does el rescate mean? (\textquotedblleft{}The rescue; the ransom\textquotedblright{}.) Say it once more.
\end{tcolorbox}
\section[salvar]{salvar — to save (a life)}

\label{lesson:ES-C526-salvar}



\begin{quote}
Before the new one: say the Spanish for to rescue, then the Spanish for the rescue.
\end{quote}

\subsection*{You'll want to know: salvar}
\textbf{salvar} — \textquotedblleft{}to save (a life)\textquotedblright{}.

\begin{quote}
El médico le salvó la vida.
\end{quote}

\begin{quote}
\emph{The doctor saved his life.}
\end{quote}

To save money is \emph{ahorrar}, not \emph{salvar}.

\begin{cousinweb}
\emph{salvar} is inherited from Late Latin \emph{salvāre} 'to save' (EtymDB), built on \emph{salvus} 'safe, unharmed'.
\end{cousinweb}

\begin{grammarlens}[title={What you've built}]
One more word for telling, explaining and describing at B1.
\end{grammarlens}

\subsection*{Your turn}
\begin{itemize}
\raggedright
  \item \emph{Say it:} \emph{salvar}
  \item \emph{Say it:} \emph{salvar}, once more
  \item \emph{Say it:} a sentence of your own with \emph{salvar}
  \item \emph{Recall:} say the Spanish for to rescue, then the Spanish for the rescue, then say \emph{salvar} again
\end{itemize}

\begin{tcolorbox}[breakable,colback=teal!4,colframe=teal!35!black,title={Before you move on}]
What does salvar mean? (\textquotedblleft{}To save (a life)\textquotedblright{}.) Say it once more.
\end{tcolorbox}
\section[sobrevivir]{sobrevivir — to survive}

\label{lesson:ES-C526-sobrevivir}



\begin{quote}
Before the new one: say the Spanish for the rescue, then the Spanish for to save.
\end{quote}

\subsection*{You'll want to know: sobrevivir}
\textbf{sobrevivir} — \textquotedblleft{}to survive\textquotedblright{}.

\begin{quote}
Todos sobrevivieron al accidente.
\end{quote}

\begin{quote}
\emph{Everyone survived the accident.}
\end{quote}

\emph{Sobrevivir a algo} — to survive something.

\begin{cousinweb}
\emph{sobrevivir} comes from Latin \emph{supervīvere} 'to outlive' (EtymDB), built on \emph{super} 'over, beyond' and \emph{vīvere} 'to live'.
\end{cousinweb}

\begin{grammarlens}[title={What you've built}]
One more word for telling, explaining and describing at B1.
\end{grammarlens}

\subsection*{Your turn}
\begin{itemize}
\raggedright
  \item \emph{Say it:} \emph{sobrevivir}
  \item \emph{Say it:} \emph{sobrevivir}, once more
  \item \emph{Say it:} a sentence of your own with \emph{sobrevivir}
  \item \emph{Recall:} say the Spanish for the rescue, then the Spanish for to save, then say \emph{sobrevivir} again
\end{itemize}

\begin{tcolorbox}[breakable,colback=teal!4,colframe=teal!35!black,title={Before you move on}]
What does sobrevivir mean? (\textquotedblleft{}To survive\textquotedblright{}.) Say it once more.
\end{tcolorbox}
\section[atrapar]{atrapar — to trap, to catch}

\label{lesson:ES-C526-atrapar}



\begin{quote}
Before the new one: say the Spanish for to save, then the Spanish for to survive.
\end{quote}

\subsection*{You'll want to know: atrapar}
\textbf{atrapar} — \textquotedblleft{}to trap, to catch\textquotedblright{}.

\begin{quote}
El humo atrapó a la familia en el piso.
\end{quote}

\begin{quote}
\emph{The smoke trapped the family in the flat.}
\end{quote}

\emph{Quedarse atrapado} — to be trapped.

\begin{cousinweb}
\emph{atrapar} is usually traced to French \emph{attraper} 'to catch in a trap', built on \emph{trappe} 'a trap', a word of Frankish origin. English \emph{trap} comes from the same Germanic family.
\end{cousinweb}

\begin{grammarlens}[title={What you've built}]
One more word for telling, explaining and describing at B1.
\end{grammarlens}

\subsection*{Your turn}
\begin{itemize}
\raggedright
  \item \emph{Say it:} \emph{atrapar}
  \item \emph{Say it:} \emph{atrapar}, once more
  \item \emph{Say it:} a sentence of your own with \emph{atrapar}
  \item \emph{Recall:} say the Spanish for to save, then the Spanish for to survive, then say \emph{atrapar} again
\end{itemize}

\begin{tcolorbox}[breakable,colback=teal!4,colframe=teal!35!black,title={Before you move on}]
What does atrapar mean? (\textquotedblleft{}To trap, to catch\textquotedblright{}.) Say it once more.
\end{tcolorbox}
